% fig_fo1_1_Uj_of_T — 평형 중심 U_j(T) 의 온도 선형 이동 (Chapter 1 §3, eq:Uj)
% 좌표 = 식 U_j(T)=(-ΔH_rxn+T ΔS_rxn)/F 의 수치 평가 (GRAPHITE_STAGING_LIT, F=96485).
% 대상 라이브러리(ch1): positioning,arrows.meta,calc,fit,backgrounds.
\begin{figure}[t]
\centering
% x = (T-258.15) K, y = (U-70) mV. 눈금: T=258/278/298/318/338, U=80/120/160/200 mV.
\begin{tikzpicture}[x=0.084cm,y=0.0265cm,>=Latex,font=\scriptsize]
% ---- T_ref 강조 띠 + 세로선 ----
\begin{scope}[on background layer]
  \fill[black!7] (37,0) rectangle (43,166);
\end{scope}
\draw[densely dashed,black!55] (40,0) -- (40,160);
\node[black!55,anchor=south,font=\tiny] at (40,160) {$T_\mathrm{ref}{=}298.15$ K};
% ---- 축 ----
\draw[->] (-3,0) -- (86,0) node[right] {$T$ [K]};
\draw[->] (-3,0) -- (-3,168) node[above] {$U_j$ [mV]};
\foreach \tx/\tl in {0/258,20/278,40/298,60/318,80/338}
  {\draw (\tx,2)--(\tx,-2) node[below] {\tl};}
\foreach \uy/\ul in {10/80,50/120,90/160,130/200}
  {\draw (-3,\uy)--(-6,\uy) node[left] {\ul};}
% ---- 네 전이 직선 U_j(T) (선형; 양끝점만으로 정확) ----
% 4->3  dS=+29  slope=+0.301 mV/K
\draw[very thick] (0,128.85) -- (80,152.90);
\fill (40,140.88) circle (1.6pt);
\node[anchor=south west,inner sep=1pt] at (63,148.2) {$4\!\to\!3$};
% 3->2L dS=0  slope 0 (수평)
\draw[thick,densely dashed] (0,69.92) -- (80,69.92);
\fill (40,69.92) circle (1.6pt);
\node[anchor=south west,inner sep=1pt] at (63,71.5) {$3\!\to\!2\mathrm L$};
% 2L->2 dS=-5  slope=-0.052 mV/K
\draw[thick,black!55] (0,52.40) -- (80,48.25);
\fill[black!55] (40,50.32) circle (1.6pt);
\node[anchor=north west,inner sep=1pt,black!55] at (63,47.5) {$2\mathrm L\!\to\!2$};
% 2->1  dS=-16  slope=-0.166 mV/K
\draw[very thick] (0,21.93) -- (80,8.66);
\fill (40,15.29) circle (1.6pt);
\node[anchor=north west,inner sep=1pt] at (63,12.0) {$2\!\to\!1$};
% ---- 기울기 = ΔS_j/F 주석 ----
\node[anchor=west,align=left,font=\tiny] at (2,150)
  {기울기 $\partial U_j/\partial T=\Delta S_{\rxn,j}/F$:\\
   $4\!\to\!3$: $+0.301$ mV/K ($\Delta S{=}{+}29$)\\
   $3\!\to\!2\mathrm L$: $0$ ($\Delta S{=}0$)};
\node[anchor=west,align=left,font=\tiny] at (2,30)
  {$2\mathrm L\!\to\!2$: $-0.052$ ($\Delta S{=}{-}5$)\\
   $2\!\to\!1$: $-0.166$ mV/K ($\Delta S{=}{-}16$)};
% T_ref 값 읽기
\node[anchor=south,font=\tiny] at (40,140.88) {};
\end{tikzpicture}
\caption{평형 중심 $U_j(T)=(-\Delta H_{\rxn,j}+T\,\Delta S_{\rxn,j})/F$(식~\eqref{eq:Uj})의 온도
의존 --- 좌표는 흑연 staging 네 전이 표준값 $(\Delta H_{\rxn,j},\Delta S_{\rxn,j})$ 를 넣은 식 그대로의
수치 평가다. 각 직선의 기울기는 정확히 $\partial U_j/\partial T=\Delta S_{\rxn,j}/F$ 라, 반응 엔트로피의
부호가 중심의 이동 방향을 가른다 --- $4\!\to\!3$($\Delta S{=}{+}29$)은 온도와 함께 오르고, $3\!\to\!2\mathrm L$
($\Delta S{=}0$)은 수평이며, $2\mathrm L\!\to\!2\cdot2\!\to\!1$($\Delta S{<}0$)은 내린다. 점은 $T_\mathrm{ref}{=}298.15$
K 의 값 $U_j{=}210.9/139.9/120.3/85.3$ mV 이고, $80$ K 창에서 중심이 수 $\sim$수십 mV 규모로 벌어진다 ---
다온도 피팅이 온도마다 $U_j(T)$ 를 따로 읽어야 하는 까닭이다.}
\label{fig:cand-fo1-1}
\end{figure}
